\input{../../../stock/en-tete_v6.tex}

\newcommand{\titrechapitre}{Théorie de la complexité}
\newcommand{\numerochapitre}{1}

\renewcommand{\headrulewidth}{0.4pt}
\renewcommand{\footrulewidth}{0.4pt}
\fancyfoot[L]{Notes de Raphaël JONTEF}
\fancyfoot[C]{\thepage}
\fancyfoot[R]{Enseirb-Matmeca}
\fancyhead[L]{Cours d'algorithmique}
\fancyhead[R]{\titrechapitre}

\begin{document}
\input{../../../stock/commands.tex}

% Page de titre custom
\setlength{\headheight}{15pt}
\begin{titlepage}
    \centering
    \vspace*{3cm}
    {\huge Chapitre \numerochapitre\par}
    {\Huge \bfseries\titrechapitre\par}
    \vspace{2cm}
    {\Large Notes rédigées par Raphaël JONTEF\par}
    \vspace{0.5cm}
    {\large Suggestions concernant le cours de Denis LAPOIRE}

    \vspace{4cm}
    {\small Enseirb-Matmeca\par}
    {\small filière informatique\par}
    \vfill
    {\large \today\par}
\end{titlepage}

% plan du cours
\tableofcontents
\newpage


% -----------Corps du document--------------

\begin{definition}{}{domination asymptotique}
    Soit $(u_n)_{n \in \N}$ et $(v_n)_{n \in \N}$ deux suites réelles positives. On dit que \notion{$(u_n)_{n \in \N}$ est dominée par $v_n$} si~:
    $$\exists C > 0,\, \forall n \in \N,\, u_n \leq C v_n$$
    On note alors $\displaystyle u_n \symbolelim{n \to + \infty}{}{=} \mc{O}(v_n)$.
\end{definition}

\begin{definition}{}{équivalence asymptotique}
    Soit $(u_n)_{n \in \N}$ et $(v_n)_{n \in \N}$ deux suites réelles positives. On dit que \notion{$(u_n)_{n \in \N}$ est équivalente à $v_n$} si~:
    $$\exists C_1 > 0,\, \exists C_2 > 0,\, \forall n \in \N,\, C_1 v_n \leq u_n \leq C_2 v_n$$
    On note alors $\displaystyle u_n \symbolelim{n \to + \infty}{}{\sim} v_n$, ou bien dans le cadre de ce cours~: $\displaystyle u_n \symbolelim{n \to + \infty}{}{=} \Theta(v_n)$.
\end{definition}

\begin{theoreme}{}{lemme de sommation}
    Soit $g : \R_+ \to \R_+$ une fonction continue et intégrable vérifiant~:
    \begin{enumeratebf}
        \item $g$ est croissante
        \item toute primitive $G$ de $g$ vérifie $G(n) \symbolelim{n \to + \infty}{}{\sim} G(n+1)$
    \end{enumeratebf}
    Alors les suites solutions de $u_n = g(n) + u_{n-1}$ sont en $\Theta\Big(G(n)\Big)$ avec $G$ une primitive de $g$.
\end{theoreme}

\begin{remarque}{}{concernant l'énoncé}
    Les conditions sont vérifiées dans le cas où $g$ est polynomiale, exponentielle ou bien encore logarithmique, et ce seront les seuls cas que nous étudierons dans un contexte d'étude de complexité.
    En particulier, supposer que $g = \mc{O}(K^n)$ comme le fait M. LAPOIRE donne l'hypothèse 2, tandis que la croissance de $g$ est assurée par sa nature de complexité.
    
\end{remarque}

\begin{demonstration}
    Soit $(u_n)_{n \in \N}$ solution. On remarque déjà que~:
    $$u_n = g(n) + g(n-1) + u_{n-2} = \sum_{k=1}^{n} g(k) + u_0$$
    Montrons alors qu'il existe $C_1 > 0$ et $C_2 > 0$ tels que~:
    $$C_1 G(n) \leq u_n \leq C_2 G(n)$$
    Soit $i \in \N^*$ et $x \in \seg{i}{i+1}$. Comme $g$ est croissante, on a~:
    $$g(x-1) \leq g(i) \leq g(x)$$
    En intégrant entre $i$ et $i+1$, on obtient par croissance de l'intégrale~:
    $$\int_{i}^{i+1} g(x-1) \d{x} \leq g(i) \leq \int_{i}^{i+1} g(x) \d{x}$$
    En sommant entre $i=1$ et $i=n$, on a par relation de Chasles~:
    $$\int_{0}^{n} g(x) \d{x} \leq \sum_{i=1}^{n} g(i) \leq \int_{1}^{n+1} g(x) \d{x}$$
    Soit~:
    $$ G(n) - G(0) \leq u_n - u_{0} \leq G(n+1) - G(1)$$
    Mais encore~:
    $$ \frac{G(n) - G(0) + u_{0}}{G(n)}G(n) \leq u_{n} \leq \frac{G(n+1) - G(1) + u_{0}}{G(n)}G(n)$$
    Ces fractions admettent pour limite 1 lorsque $n$ tend vers $+\infty$ (par hypothèse). Elles sont donc bornées à partir d'un certain rang et en déduit l'existence de $C_1$ et $C_2$.
\end{demonstration}

\end{document}